\section{Discussion and limitations}\label{sec:discussion}

The acquisition region in \cref{obj:def:annotation-region} provides a benchmark for allocating response labels and routine covariate-treatment records under the model in \cref{obj:def:primitive-class,obj:ass:sample-law}. It identifies sample-size pairs for which the sharp rate is within a specified factor of the supplied-propensity order. Through \cref{obj:thm:sharp-annotation-frontier}, this rate comparison translates into a minimax precision comparison up to constants depending on the public class parameters. For an absolute-error target, \cref{obj:prop:annotation-threshold} instead gives simultaneous requirements on the labeled count and the product of labeled and total treatment-record counts. These requirements offer a basis for comparing acquisition plans once their costs are specified.

A chart-review setting illustrates the distinction between the two record types. A labeled record contains covariates, treatment, and a true response ascertained through review; an outcome-free record contains the covariates and treatment available from routine records. Adding a new independent labeled record changes \((n,m)\) to \((n+1,m)\) and increases the total count. Reviewing a response for one record already present in the auxiliary block changes the composition to \((n+1,m-1)\) at fixed total count; if the record index is fixed in advance, or chosen by a randomization independent of all records, the resulting dataset follows the stated sampling law at \((n+1,m-1)\). Adding a new routine record changes \((n,m)\) to \((n,m+1)\). The sample requirements benchmark these alternative compositions under the specified independent, common-population experiment; acquisition costs determine which transition is relevant in an application.

The two channels address distinct sources of statistical difficulty. Response labels carry information about the probability-valued conditional response means and the causal contrast. The arm-mean conditions in \cref{obj:ass:control-interior,obj:ass:treated-interior} place both means in \([0,1]\), and \cref{obj:thm:supplied-propensity} characterizes the labeled-sample precision when the propensity function is supplied. Outcome-free records carry information about treatment assignment conditional on covariates. Their contribution enters the second component of the sharp rate in \cref{obj:def:sharp-rate}, where the product of labeled and total counts governs the cost of nuisance adjustment. Increasing the routine-record count can reduce that component until the supplied-propensity order determines the overall rate.

Public nuisance regularity determines how much routine information is needed to reach this benchmark. The rate exponent depends on the sum of the propensity and control-mean H\"older orders, while the effect H\"older order governs the supplied-propensity precision. Greater nuisance regularity can reduce the total-record growth needed at a given labeled count. In the sufficiently smooth regime characterized by \cref{obj:thm:sharp-annotation-frontier}, labeled records alone attain the supplied-propensity order; in the lower-regularity regime, additional outcome-free records can improve the minimax order until saturation. These comparisons concern rates, with constants depending on the fixed class parameters. Because the regularity parameters are public, they also determine the estimator's localization and projection choices in \cref{obj:def:tuning}. Acquisition planning and estimator tuning thus use the same description of the statistical model.

\subsection{Limitations and future work}

The analysis uses a known uniform covariate distribution on the unit cube, a fixed public overlap level \(\varepsilon\in(0,1/2)\), binary responses with arm means in \([0,1]\), and independent record blocks from a common population. Extending the acquisition benchmark to general, potentially unknown covariate distributions is a direction for future work. Covariate-dependent or adaptive selection of records for review would also change the labeled-record law and requires a sampling or missingness model beyond the experiment studied here. Covariate shift between reviewed and routine records would introduce an additional transport problem. These extensions would help connect the benchmark to populations with uneven support and heterogeneous record availability.

Unknown smoothness raises a separate question of adaptation. The present tuning and acquisition comparisons use public H\"older orders and a public radius. Procedures that select tuning parameters from the data would need risk guarantees over multiple regularity classes, together with an assessment of any adaptation cost. Uncertainty quantification also requires further analysis: the expected absolute-error bounds characterize point-estimation precision, whereas confidence intervals require coverage guarantees and control of approximation bias.

Alternative losses and causal targets may change the allocation requirements. Squared-error loss, integrated estimation of the treatment-effect function, average treatment effects, and policy-value targets each call for their own precision benchmark. Richer auxiliary information is another direction, including response proxies, repeated measurements, and additional variables predictive of the response. Such information would enlarge the observation model beyond covariate-treatment pairs and require an explicit account of how the auxiliary measurements relate to the true response.
